%══════════════════════════════════════════════════════════════════════════════
% MAIN-STORY.TEX — Template para historias científicas verticales 9:16
% Motor: pdfLaTeX  |  Compilar DOS veces (TikZ remember picture)
%
% ZONA SEGURA para redes sociales:
%   Columnas:  2-11 (10 columnas de contenido; 1 y 12 = márgenes)
%   Filas:     3-22 (20 filas de contenido; 1-2 y 23-24 = buffers)
%
% LAYOUT COMPLETO (24 filas × 12 columnas):
%   Fila  1-2    → Metadatos, acento, categoría (sangran OK)
%   Filas 3-6    → HÉROE: ecuación (SAFE)
%   Fila  7      → Barra de acento (SAFE)
%   Filas 8-10  → Título (SAFE)
%   Fila  11     → Subtítulo (SAFE)
%   Filas 12-14  → Abstract (SAFE)
%   Filas 15-22  → Figura/Opción 1 (SAFE)
%   Fila  15-22  → Body/Opción 2 (SAFE)
%   Filas 23-24  → Footer (sangra OK)
%══════════════════════════════════════════════════════════════════════════════
\documentclass[9pt,oneside]{article}

%── 1. Paquetes ───────────────────────────────────────────────────────────────
%── config/packages-story.tex ──────────────────────────────────────────────
% Subconjunto mínimo para póster de una sola página.
%───────────────────────────────────────────────────────────────────────────

% ─── Encoding y lenguaje ────────────────────────────────────────────────────
\usepackage[utf8]{inputenc}
\usepackage[T1]{fontenc}
\usepackage[spanish,es-noshorthands]{babel}
\deactivatequoting
\usepackage{csquotes}

% ─── Tipografía ─────────────────────────────────────────────────────────────
% Inconsolata: monospace con carácter, garantizado en TeX Live estándar.
% \familydefault → \ttdefault hace que TODO el texto sea monospace.
% Las matemáticas conservan Computer Modern (fuente propia de math mode).
\usepackage{inconsolata}
\renewcommand{\familydefault}{\ttdefault}
\usepackage{microtype}

% ─── Matemáticas ────────────────────────────────────────────────────────────
\usepackage{amsmath}
\usepackage{amssymb}
\usepackage{mathtools}
\usepackage{bm}
\usepackage{physics}
\usepackage{cancel}
\usepackage{siunitx}
\sisetup{
  output-decimal-marker = {,},
  per-mode              = symbol,
}

% ─── Color ──────────────────────────────────────────────────────────────────
\usepackage{xcolor}

% ─── Cajas y poster ─────────────────────────────────────────────────────────
\usepackage[many,poster]{tcolorbox}
\tcbuselibrary{skins,hooks,poster,fitting}

% ─── TikZ ───────────────────────────────────────────────────────────────────
\usepackage{tikz}
\usetikzlibrary{
  arrows.meta,
  positioning,
  calc,
  decorations.pathreplacing,
  backgrounds,
}
\usepackage{pgfplots}
\pgfplotsset{compat=1.18}

% ─── Figuras ────────────────────────────────────────────────────────────────
\usepackage{graphicx}

% ─── Listings ───────────────────────────────────────────────────────────────
\usepackage{listings}

% ─── Utilidades ─────────────────────────────────────────────────────────────
\usepackage{etoolbox}
\usepackage{xparse}
\usepackage[hidelinks]{hyperref}

%── 2. Geometría ───────────────────────────────────────────────────────────────
%── config/geometry-story.tex ──────────────────────────────────────────────
% Página vertical 9:16 escalada a 13.5×24 cm para mayor resolución
% tipográfica. Ratio exacto 9:16 = 0.5625.
% Sin headers, footers ni márgenes generosos: todo el espacio es poster.
%───────────────────────────────────────────────────────────────────────────
\usepackage[
  paperwidth  = 13.5cm,
  paperheight = 24.0cm,
  top         = 0.35cm,
  bottom      = 0.35cm,
  left        = 0.35cm,
  right       = 0.35cm,
  noheadfoot,
  nomarginpar,
]{geometry}

\pagestyle{empty}
\setlength{\parindent}{0pt}
\setlength{\parskip}{0pt}

% ─── Eliminar espacios verticales por defecto del article class ──────────────
\setlength{\topmargin}{-1in}
\setlength{\headheight}{0pt}
\setlength{\headsep}{0pt}
\setlength{\footskip}{0pt}
\setlength{\textheight}{\paperheight}
\raggedbottom

%── 3. Paleta de colores ───────────────────────────────────────────────────────
\input{themes/catppuccin-palette.tex}
%── themes/story-accents.tex ────────────────────────────────────────────────
%
% Paleta tonal del poster: grises con tinte azul-índigo.
% Son los overlays de Catppuccin Mocha usados sobre blanco —
% más sofisticados que grises puros, coherentes con el sistema Ctp.
%
% Sistema de acento: un único color por historia, definido en
% story-content.tex como \StoryAccentColor (nombre de color Ctp*).
% Se materializa en \begin{document} con:
%   \colorlet{StoryAccent}{\StoryAccentColor}
%   \colorlet{StoryAccentFaint}{\StoryAccentColor!12!white}
%───────────────────────────────────────────────────────────────────────────

% ─── Escala tonal (4 niveles + fantasma) ────────────────────────────────────
\definecolor{InkBlack}{HTML}{1e1e2e}   % héroe, énfasis máximo
\definecolor{InkDark} {HTML}{313244}   % títulos principales
\definecolor{InkMid}  {HTML}{585b70}   % body text, subtítulos
\definecolor{InkLight}{HTML}{9399b2}   % metadatos, captions, fechas
\definecolor{InkFaint}{HTML}{cdd6f4}   % líneas sutiles, número decorativo

% ─── Nota: NO se sobreescriben colores Ctp* aquí ────────────────────────────
% La paleta Catppuccin completa sigue disponible para ecuaciones
% (equation-styles.tex) y para el acento de color de la historia.

%── 4. Macros matemáticos ──────────────────────────────────────────────────────
\input{config/macros.tex}

%── 5. Configuraciones que requieren colores ────────────────────────────────────
\input{config/equation-styles.tex}
\input{config/pgfplots-config.tex}

%── 6. Tipografía y cajas ──────────────────────────────────────────────────────
%── config/typography-story.tex ─────────────────────────────────────────────
%
% Jerarquía tipográfica del poster.
% Contraste de escala: héroe 16pt ↔ metadatos 5pt (ratio 3.2×).
%───────────────────────────────────────────────────────────────────────────

% ─── Número decorativo (40pt, InkFaint) ─────────────────────────────────────
\newcommand{\DisplayNum}[1]{%
  {\fontsize{40}{40}\selectfont\bfseries\color{InkFaint}#1}%
}

% ─── Ecuación héroe (escala CM math vía fontsize previa al entorno math) ─────
\newcommand{\HeroFont}{\fontsize{16}{20}\selectfont}

% ─── Título ──────────────────────────────────────────────────────────────────
\newcommand{\TitleFont}{\fontsize{30}{23}\selectfont\bfseries\color{InkBlack}}

% ─── Subtítulo ───────────────────────────────────────────────────────────────
\newcommand{\SubFont}{\fontsize{9}{11}\selectfont\color{InkMid}}

% ─── Body y abstract ─────────────────────────────────────────────────────────
\newcommand{\BodyFont}{\fontsize{7.5}{10.5}\selectfont\color{InkMid}}

% ─── Caption de figura ───────────────────────────────────────────────────────
\newcommand{\CaptFont}{\fontsize{6.5}{8}\selectfont\color{InkLight}}

% ─── Chip de categoría ───────────────────────────────────────────────────────
\newcommand{\CatFont}{\fontsize{6}{7}\selectfont\bfseries}

% ─── Metadatos técnicos (barras superior e inferior) ─────────────────────────
\newcommand{\MetaFont}{\fontsize{8}{8}\selectfont\color{InkLight}}

% ─── Footer (referencias y marca) ────────────────────────────────────────────
\newcommand{\FootFont}{\fontsize{5.5}{7}\selectfont\color{InkLight}}
%── config/tcolorbox-story.tex ──────────────────────────────────────────────
%
% Filosofía: ninguna caja es visible.
% tcolorbox actúa exclusivamente como sistema de layout (posicionamiento,
% padding, alineación). El diseño visible emerge solo de tipografía,
% espacio y los dos únicos elementos cromáticos: \AccentBar y \StoryCat.
%───────────────────────────────────────────────────────────────────────────

% ─── Estilos base ────────────────────────────────────────────────────────────

% ghost: caja completamente invisible, sin padding
\tcbset{
  ghost/.style={
    enhanced,
    frame hidden,
    colback=white,
    boxrule=0pt,
    left=0pt, right=0pt, top=0pt, bottom=0pt,
    sharp corners,
  },
}

% ghostpad: ghost con padding mínimo (para separar contenido del borde del grid)
\tcbset{
  ghostpad/.style={
    ghost,
    left=4pt, right=4pt, top=3pt, bottom=3pt,
  },
}

% leftaccent: borde izquierdo en color de acento, sin fondo visible
% Uso: [leftaccent=StoryAccent]
\tcbset{
  leftaccent/.style={
    ghost,
    borderline west={1.5pt}{0pt}{#1},
    left=8pt,
    top=2pt, bottom=2pt, right=4pt,
  },
}

% ─── Elementos cromáticos (los únicos dos elementos de color del poster) ─────

% Chip de categoría: único bloque con fondo tintado
% Uso: \StoryCat{StoryAccent}{Nombre de Categoría}
\newcommand{\StoryCat}[2]{%
  \tcbox[
    enhanced,
    frame hidden,
    colback=#1!12!white,
    arc=1.5pt,
    boxrule=0pt,
    left=4pt, right=4pt, top=1.5pt, bottom=1.5pt,
    nobeforeafter,
    box align=base,
  ]{\CatFont\color{#1}\MakeUppercase{#2}}%
}

% Barra horizontal de acento: 2pt de color puro, ancho completo
% Es el único elemento de color sólido en el poster (referencia: los
% rectángulos naranja/negro de la imagen japonesa).
% Uso: \AccentBar{StoryAccent}
\newcommand{\AccentBar}[1]{%
  {\color{#1}\rule{\linewidth}{2pt}}%
}

% ─── Separadores ─────────────────────────────────────────────────────────────

% Línea gris sutil (footer, separaciones internas)
\newcommand{\GrayRule}{%
  {\color{InkFaint}\rule{\linewidth}{0.4pt}}%
}

%── 7. Entornos ───────────────────────────────────────────────────────────────
\input{environments/generator.tex}
%── environments/story-blocks.tex ───────────────────────────────────────────
%
% Entornos auxiliares del poster.
% IMPORTANTE: todos usan StoryAccent y StoryAccentFaint (nombres resueltos
% con \colorlet en main-story.tex), nunca \StoryAccentColor directamente.
%───────────────────────────────────────────────────────────────────────────

% ─── Highlight inline de expresión matemática ────────────────────────────────
% Fondo muy sutil en el color de acento de la historia.
% Uso: la función \inlinemath{f(x;\theta)} parametriza la densidad.
\newcommand{\inlinemath}[1]{%
  \tcbox[
    enhanced,
    frame hidden,
    colback=StoryAccentFaint,
    arc=1.5pt,
    boxrule=0pt,
    left=1.5pt, right=1.5pt, top=0pt, bottom=0pt,
    nobeforeafter,
    box align=base,
  ]{\small$\displaystyle#1$}%
}

% ─── Dato clave ──────────────────────────────────────────────────────────────
% Bloque compacto con borde izquierdo, sin fondo.
% Uso: \KeyFact{La varianza estacionaria es $\sigma^2/(2\theta)$}
\newcommand{\KeyFact}[1]{%
  \begin{tcolorbox}[
    leftaccent=StoryAccent,
    fontupper=\fontsize{7}{9}\selectfont\bfseries\color{InkBlack},
  ]
    #1
  \end{tcolorbox}%
}

% ─── Separador con label opcional ────────────────────────────────────────────
% Uso: \StorySep  o  \StorySep{Derivación}
\newcommand{\StorySep}[1][]{%
  \par\vspace{3pt}%
  \GrayRule%
  \ifx&#1&%
  \else
    \par\vspace{1pt}%
    {\fontsize{5.5}{6}\selectfont\bfseries\color{InkLight}\MakeUppercase{#1}}%
  \fi
  \par\vspace{3pt}%
}

%── 8. Contenido de la historia ────────────────────────────────────────────────
%══════════════════════════════════════════════════════════════════════════════
% STORY-CONTENT.TEX — Historia #001: Score Matching y Difusión Estocástica
%
% Único archivo a editar por historia.
% \providecommand (en lugar de \newcommand) evita errores si el archivo
% se carga más de una vez accidentalmente.
%══════════════════════════════════════════════════════════════════════════════

% ─── Identidad ───────────────────────────────────────────────────────────────
\providecommand{\StoryAccentColor}{CtpMauve}
\providecommand{\StoryNumber}{001}
\providecommand{\StoryDate}{2025}
\providecommand{\StorySeriesTag}{QbitLab Stories}
\providecommand{\StoryCategory}{Modelos Generativos}

% ─── Textos ──────────────────────────────────────────────────────────────────
\providecommand{\StoryTitle}{%
  \textquestiondown C\'OMO APRENDE\\ UNA RED A\\ GENERAR RUIDO?%
}
\providecommand{\StorySubtitle}{\textbf{Score matching y difusi\'on estoc\'astica}}
\providecommand{\StoryModelLabel}{Objetivo de entrenamiento (denoising score matching)}

\providecommand{\StoryAbstract}{%
  Los modelos de difusi\'on aprenden a invertir un proceso de corrupci\'on
  gaussiana. El truco central: en lugar de estimar $p(x)$ directamente
  ---intractable--- se estima su gradiente logar\'itmico, el \emph{score}
  $\nabla_x \log p(x)$. Gracias a la identidad de Vincent (2011), esto
  equivale a predecir el ruido a\~nadido: regresi\'on pura.%
}

% ─── Ecuación héroe ──────────────────────────────────────────────────────────
% Renderizada a ~16pt en display mode. Usa macros de equation-styles.tex.
\providecommand{\StoryEquation}{%
  \eqfn{\mathcal{L}}(\equ{\theta})
  \;=\;
  \eqcon{\tfrac{1}{2}}\,
  \eqop{\mathbb{E}}_{\!\eqk{x}}
  \!\left[
    \bigl\|
      \eqfn{s_{\equ{\theta}}}(\eqk{x})
      -
      \eqop{\nabla_{\!\eqk{x}}} \log \eqk{p(x)}
    \bigr\|^2
  \right]%
}

% ─── Figura ───────────────────────────────────────────────────────────────────
\providecommand{\StoryFigure}{%
  \begin{tikzpicture}
    \begin{axis}[
      catppuccin-clean,
      width       = 10cm,
      height      = 6cm,
      xlabel      = {$t$},
      ylabel      = {$\sigma(t)$},
      title       = {Programaci\'on de ruido},
      domain      = 0:1,
      samples     = 80,
      no markers,
      legend pos  = north west,
      title style      = {font=\fontsize{7}{8}\selectfont\bfseries,
                          color=InkDark},
      label style      = {font=\fontsize{6.5}{7}\selectfont,
                          color=InkMid},
      tick label style = {font=\fontsize{6}{7}\selectfont,
                          color=InkLight},
      legend style     = {font=\fontsize{6}{7}\selectfont,
                          draw=InkFaint, fill=white, text=InkMid},
      axis line style  = {color=InkFaint},
      tick style       = {color=InkFaint},
    ]
      \addplot+[thick] {sqrt(1 - exp(-5*x))};
      \addlegendentry{VP-SDE}
      \addplot+[thick] {x};
      \addlegendentry{VE-SDE}
      \addplot+[thick, dashed] {sqrt(x*(1-x))};
      \addlegendentry{Sub-VP}
    \end{axis}
  \end{tikzpicture}%
}

\providecommand{\StoryFigCaption}{%
  Tres esquemas de programaci\'on de ruido. El score $s_\theta(x,t)$
  se entrena para cada nivel $t \in [0,1]$.%
}

% ─── Body ─────────────────────────────────────────────────────────────────────
\providecommand{\StoryBody}{%
  Durante la generaci\'on, el score guia una SDE de Langevin en tiempo
  revertido: el proceso de Ornstein--Uhlenbeck inverso. La curva VP-SDE
  preserva la varianza total; VE-SDE la hace crecer. Ambas convergen a
  la misma distribuci\'on cuando el score es exacto.%
}

% ─── Referencias ─────────────────────────────────────────────────────────────
\providecommand{\StoryRefs}{%
  Song et al.\ (2021) \textit{Score-Based Generative Modeling through SDEs.}
  ICLR.\\
  Vincent (2011) \textit{A Connection Between Score Matching
  and Denoising Autoencoders.} Neural Comput.%
}

%── 9. Metadatos PDF ───────────────────────────────────────────────────────────
\hypersetup{
  pdftitle  = {\StoryTitle},
  pdfauthor = {QbitLab},
}

%── 10. Color de página (preamble, antes de \begin{document}) ──────────────────
\pagecolor{white}
\color{InkDark}

% Materializar color de acento
\colorlet{StoryAccent}{\StoryAccentColor}
\colorlet{StoryAccentFaint}{\StoryAccentColor!12!white}

%══════════════════════════════════════════════════════════════════════════════
\begin{document}

%══════════════════════════════════════════════════════════════════════════════
% POSTER GRID  (24 filas × 12 columnas)
% Zona segura: columnas 2-11, filas 3-22
%══════════════════════════════════════════════════════════════════════════════
\vspace*{0pt}%
\begin{tcbposter}[
  poster={
    showframe = true,      % cambiar a true para ver grid (debug)
    columns   = 12,
    rows      = 24,
    spacing   = 1.2mm,
    height    = \paperheight - 0.7cm,
    width     = \paperwidth  - 0.7cm,
  },
  boxes={ghost},
]

%──────────────────────────────────────────────────────────────────────────────
% [F1-2] METADATOS + ACENTO + CATEGORÍA
% Filas 1-2 son buffers (pueden sangrar fuera de la zona segura).
% Esto permite que la categoría y metadatos sean visibles pero no críticos.
%──────────────────────────────────────────────────────────────────────────────

%── [F1] Metadatos técnicos (serie, fecha, categoría) ────────────────────────
\posterbox[ghostpad]{column=2, span=10, row=1}{%
  \MetaFont
  \begin{minipage}[c]{0.6\linewidth}
    \StorySeriesTag\enspace ---\enspace \StoryDate
  \end{minipage}%
  \hfill
  \begin{minipage}[c]{0.35\linewidth}\raggedleft
    \StoryCategory
  \end{minipage}%
}

%── [F2a] Barra de acento (fondo, todo el ancho) ───────────────────────────────
\posterbox[ghost, top=0pt, bottom=0pt, left=4pt, right=4pt, valign=center]{%
  column=2, span=8, row=2%
}{%
  \AccentBar{StoryAccent}%
}
 
%── [F2b] Chip de categoría (flotante sobre la línea, derecha) ─────────────────
\posterbox[
  ghost,
  top=0pt, bottom=0pt, left=10pt, right=0pt,
  halign=right, valign=center
]{%
  column=9, span=2, row=2%
}{%
  \StoryCat{StoryAccent}{\StoryCategory}%
}

%══════════════════════════════════════════════════════════════════════════════
% ZONA SEGURA COMIENZA (filas 3-22, columnas 2-11)
%══════════════════════════════════════════════════════════════════════════════

%──────────────────────────────────────────────────────────────────────────────
% [F3–6] HÉROE: ECUACIÓN PRINCIPAL (SAFE)
% 6 filas = ~25% del área. Centro del espacio seguro.
%──────────────────────────────────────────────────────────────────────────────
\posterbox[ghost, halign=center, valign=center]{%
  column=2, span=10, row=3, rowspan=4%
}{%
  {\HeroFont
  \begin{equation*}
    \displaystyle\StoryEquation
  \end{equation*}}%
  \par\vspace{4pt}%
  {\CaptFont\StoryModelLabel}%
}

%──────────────────────────────────────────────────────────────────────────────
% [F7] BARRA DE ACENTO (SAFE)
% Divisor visual entre sección matemática y editorial.
%──────────────────────────────────────────────────────────────────────────────
\posterbox[ghost, top=0pt, bottom=0pt, left=4pt, right=4pt, valign=center]{%
  column=2, span=10, row=7%
}{%
  \AccentBar{StoryAccent}%
  \vspace{2pt}%
}

%──────────────────────────────────────────────────────────────────────────────
% [F8–10] TÍTULO + NÚMERO DECORATIVO (SAFE)
% Título en columnas 2-11. Número decorativo en "zona de margen" (columna 12).
%──────────────────────────────────────────────────────────────────────────────
\posterbox[ghostpad]{column=2, span=9, row=8, rowspan=3}{%
  {\TitleFont\StoryTitle\par}%
}

% Número decorativo: en margen derecho (column=12, puede sangrar)
\posterbox[ghost, halign=right, valign=bottom]{%
  column=10, span=1, row=8, rowspan=1%
}{%
  \DisplayNum{\StoryNumber}%
}

%──────────────────────────────────────────────────────────────────────────────
% [F11] SUBTÍTULO (SAFE)
%──────────────────────────────────────────────────────────────────────────────
\posterbox[ghostpad]{column=2, span=9, row=11}{%
  {\SubFont\StorySubtitle}%
}

%──────────────────────────────────────────────────────────────────────────────
% [F12–14] ABSTRACT (SAFE)
% Borde izquierdo de acento: único indicio estructural visible.
%──────────────────────────────────────────────────────────────────────────────
\posterbox[leftaccent=StoryAccent]{column=2, span=9, row=12, rowspan=3}{%
  {\BodyFont\StoryAbstract}%
}

%──────────────────────────────────────────────────────────────────────────────
% [F15–22]  FIGURA / GRÁFICO (SAFE)
% 8 filas = ~33% del área.
%──────────────────────────────────────────────────────────────────────────────
\posterbox[ghost, halign=center, valign=center]{%
  column=2, span=9, row=15, rowspan=8%
}{%
  \centering
  \StoryFigure
  \par\vspace{3pt}%
  {\CaptFont\StoryFigCaption}%
}

%──────────────────────────────────────────────────────────────────────────────
% [F15–22] PÁRRAFO DESCRIPTIVO / ALTERNATIVA A FIGURA (SAFE)
%──────────────────────────────────────────────────────────────────────────────
%\posterbox[ghostpad]{column=2, span=9, row=15,  rowspan=8}{%
%  {\BodyFont\StoryBody}%
%}

%──────────────────────────────────────────────────────────────────────────────
% [F23-24] FOOTER: REFERENCIAS + MARCA (BUFFERS, pueden sangrar)
% Filas 23-24 son buffers inferiores (pueden cortarse en redes sociales).
%──────────────────────────────────────────────────────────────────────────────
\posterbox[
  ghost,
  borderline north={0.4pt}{1pt}{InkFaint},
  left=4pt, right=4pt, top=2pt, bottom=2pt,
]{%
  column=2, span=10, row=23, rowspan=2%
}{%
  \begin{minipage}[c]{0.74\linewidth}
    {\fontsize{5.5}{7}\selectfont\color{InkLight}\StoryRefs}%
  \end{minipage}%
  \hfill
  \begin{minipage}[c]{0.22\linewidth}\raggedleft
    {\fontsize{5.5}{7}\selectfont\color{StoryAccent}@QbitLab}%
  \end{minipage}%
}

\end{tcbposter}
\end{document}